\section{Konklusjon}
Målingene ved 2498, 3998 og 6001 rpm viste at den største vibrasjonen lå ved grunnfrekvensen. Økningen var særlig tydelig ved 6001 rpm, der amplituden nådde 2{,}239\,mm/s. Dette er forenlig med ubalanse eller andre rotasjonsrelaterte krefter, og er forventet for et normalt fungerende system.

Amplitudene ved lagerfrekvensene $f_y$, $f_i$ og $f_r$ var svært lave ved alle turtall. Forsøket gir derfor ingen tydelige indikasjoner på lokal skade i ytterring, innerring eller kule. De lave amplitudene tilsier tilstandsgrad 1 med hensyn til lokal lagerskade, og det anbefales å videreføre normal drift med periodisk vibrasjonsovervåking.
